\section{Appendix: Upper-bound proofs}\label{sec:upper-bound-proofs}

The upper bounds in \cref{obj:thm:rectangular-upper,obj:thm:supplied-propensity} combine local polynomial approximation with control of the corresponding empirical moments. For the rectangular estimator, the two-scale construction in \cref{obj:def:projection} and the public role split in \cref{obj:def:roles} are statistical inputs: the coarse polynomial targets the smooth contrast, the fine projection controls nuisance approximation, and the independent roles separate response information from treatment-assignment information. The auxiliary results below follow this dependency order.

\subsection{Local approximation and projection}

Fix a primitive law \(P\) from \cref{obj:def:primitive-class} and admissible localization and resolution inputs \(h,J\). Retain the notation \(\bar R\) for the expected response vector and \(Q_{P,h,J}\) for the expected symmetrized matrix:
\[
\bar R=\mathbb E_{\mathsf E_{n,m}(P)}\widehat R_{h,J},
\qquad
Q_{P,h,J}=\mathbb E_{\mathsf E_{n,m}(P)}\widehat Q_{h,J}.
\]
Both expectations concern the empirical moments in \cref{obj:def:moments}. Let \(p\) denote the Taylor polynomial of \(\tau_P\) at \(x_0\), with degree \(p_*=\lceil\gamma\rceil-1\), and let \(\vartheta\) denote its coefficient vector in the coarse basis. Since \(1\le\gamma\le3\), this Taylor degree lies within the degree-two polynomial space used by the estimator. The following result collects the approximation and projection properties needed for the population estimating equations and their sampling bounds.

\begin{lemmav}{lem:local-polynomial-projection-facts}[Local polynomial projection bounds]
Fix public parameters \(d\in\mathbb N\) and \(\alpha,\beta,\gamma,L,\varepsilon\in\mathbb R\) satisfying:
\begin{itemize}
\item \textbf{(Dimension.)} \(d\ge1\).
\item \textbf{(Smoothness.)} \(0<\alpha\le1\), \(0<\beta\le1\), and \(1\le\gamma\le3\).
\item \textbf{(Radius.)} \(L>1/2\).
\item \textbf{(Overlap.)} \(0<\varepsilon<1/2\).
\end{itemize}
There exists a finite constant \(C=C(d,\alpha,\beta,\gamma,L,\varepsilon)>0\) such that, for every primitive law \(P\in\mathcal M(d,\alpha,\beta,\gamma,L,\varepsilon)\) as defined in \cref{obj:def:primitive-class}, every \(h\in\mathbb R\) with \(0<h\le1/2\), and every integer \(J\ge1\), the following conclusions hold for the bases and projection of \cref{obj:def:projection}.

Write \(x_0=(1/2,\ldots,1/2)\) for the evaluation center, \(C_h=[1/2-h/2,1/2+h/2]^d\) for the localization cube, and \(\nu_h\) for its uniform probability measure. Let \(p\) be the Taylor polynomial of the canonical contrast \(\tau_P\) at \(x_0\), with Taylor degree \(p_*=\lceil\gamma\rceil-1\):
\[
p(x)=
\sum_{\substack{\kappa\in\mathbb N^d\\ \sum_{i=1}^d\kappa_i\le p_*}}
\frac{\partial^\kappa\tau_P(x_0)}
     {\prod_{i=1}^d\kappa_i!}
\prod_{i=1}^d(x_i-1/2)^{\kappa_i}.
\]
There exists a coarse coefficient vector \(\vartheta\) whose represented polynomial
\[
p_\vartheta(x)=r(x)^\top\vartheta
\]
satisfies
\[
\sup_{x\in C_h}|\tau_P(x)-p_\vartheta(x)|\le Ch^\gamma,
\qquad
\|\vartheta\|_2\le C,
\qquad
p_\vartheta(x_0)=\tau_P(x_0),
\]
and
\[
p_\vartheta(x)=p(x)\qquad\text{for every }x\in C_h.
\]

The coarse basis vector \(r\) and fine projection kernel \(K_J\) satisfy
\[
\sup_{x\in C_h}\|r(x)\|_2\le C,
\qquad
\sup_{x\in C_h}\int |K_J(x,x')|\,d\nu_h(x')\le C,
\]
and
\[
\iint K_J(x,x')^2\,d\nu_h(x)\,d\nu_h(x')
=qJ^d,
\]
where \(q\) is the number of coarse multi-indices of total degree at most two:
\[
q=\#\left\{\kappa\in\mathbb N^d:\sum_{i=1}^d\kappa_i\le2\right\}.
\]

Finally, with fine-cell side length \(\delta=h/J\), every coarse basis coordinate \(r_u\) satisfies
\[
\|e_Pr_u-\Pi_J(e_Pr_u)\|_{L_2(\nu_h)}
\le C\delta^\alpha,
\]
where \(e_P\) is the designated admissible propensity witness, and the designated admissible control-mean witness \(\mu_{0,P}\) satisfies
\[
\|\mu_{0,P}-\Pi_J\mu_{0,P}\|_{L_2(\nu_h)}
\le C\delta^\beta.
\]
\end{lemmav}

The Taylor approximation controls variation of the contrast across the localization cube while reproducing its value exactly at the evaluation profile. The bounded coefficient vector and coarse basis provide uniform control of polynomial evaluation. For the fine projection, the integral bound controls its action on bounded approximation errors, whereas the squared-kernel identity records the sampling cost of its rank. The two nuisance residual bounds concern different functions: the propensity multiplied by a coarse basis coordinate, and the control mean. Their respective orders account for the combined nuisance approximation term in \cref{obj:thm:rectangular-upper}.

\subsection{Population positivity}

The corrected empirical matrix has a population counterpart whose stability is determined by overlap and orthogonal projection geometry. The next identity separates its quadratic form into an overlap-weighted polynomial norm and a squared projection residual.

\begin{lemmav}{lem:population-positivity}[Population quadratic form positivity]
Suppose the following conditions hold:
\begin{itemize}
\item \textbf{(Public parameters.)} The dimension \(d\) is an integer with \(d\ge1\), the smoothness parameters satisfy \(0<\alpha,\beta\le1\) and \(1\le\gamma\le3\), the radius satisfies \(L>1/2\), and the overlap parameter satisfies \(0<\varepsilon<1/2\).
\item \textbf{(Primitive law.)} The law \(P\) belongs to the primitive class \(\mathcal M\) at these public parameters, as specified in \cref{obj:def:primitive-class}. Let \(e_P\) be its designated admissible propensity witness.
\item \textbf{(Sample sizes.)} The sample counts \(n,m\) are nonnegative integers with \(n\ge2\).
\item \textbf{(Localization and projection.)} The localization scale satisfies \(0<h\le1/2\), and the grid resolution \(J\) is an integer with \(J\ge1\). Use the projection \(\Pi_J\) from \cref{obj:def:projection}.
\end{itemize}
Let \(\nu_h\) be the uniform localization measure on the localization cube \(C_h\), with
\[
d\nu_h(x)=h^{-d}\mathbf 1_{C_h}(x)\,dx.
\]
For every coefficient vector \(v\in\mathbb R^q\), indexed by the multi-indices of total degree at most two, define its coarse-basis polynomial and the population matrix by
\[
p_v(x)=\sum_u r_u(x)v_u,
\qquad
Q_{P,h,J}
=\mathbb E_{\mathsf E_{n,m}(P)}[\widehat Q_{h,J}],
\]
where \(r_u\) are the scaled orthonormal coarse polynomial basis functions and \(\widehat Q_{h,J}\) is the symmetrized corrected empirical matrix. Then
\[
\begin{aligned}
v^\top Q_{P,h,J}v
&=\int e_P(x)\bigl(1-e_P(x)\bigr)p_v(x)^2\,d\nu_h(x)\\
&\quad+\int
\bigl[e_P(x)p_v(x)-\Pi_J(e_Pp_v)(x)\bigr]^2\,d\nu_h(x)\\
&\ge \varepsilon(1-\varepsilon)\sum_u v_u^2.
\end{aligned}
\]
\end{lemmav}

The overlap restriction in \cref{obj:ass:overlap} supplies the lower bound through the first integral. The second integral is nonnegative and makes the contribution of finite projection explicit. Thus the population matrix retains a fixed positivity margin at every admissible resolution. This margin supplies the population reference for the empirical eigenvalue guard in \cref{obj:def:estimator}.

\subsection{Rectangular sampling}

The public split assigns disjoint observations to the outcome and treatment roles under \cref{obj:ass:sample-law}. Their counts in \cref{obj:def:roles} are proportional to \(n\) and \(N=n+m\), respectively, including when \(m=0\). Consequently, the sampling analysis distinguishes ordinary localized averages from the rectangular averages formed across the two roles. The next result controls both the response vector and the matrix used for inversion.

\begin{lemmav}{lem:rectangular-sampling-bound}[Rectangular sampling bound]
Fix the covariate dimension \(d\in\mathbb N\), the smoothness parameters \(\alpha,\beta,\gamma\in\mathbb R\), the Hölder radius \(L\in\mathbb R\), and the overlap parameter \(\varepsilon\in\mathbb R\), satisfying
\begin{itemize}
\item \textbf{(Dimension.)} \(d\ge1\).
\item \textbf{(Nuisance smoothness.)} \(0<\alpha\le1\) and \(0<\beta\le1\).
\item \textbf{(Effect smoothness.)} \(1\le\gamma\le3\).
\item \textbf{(Radius.)} \(L>1/2\).
\item \textbf{(Overlap parameter.)} \(0<\varepsilon<1/2\).
\end{itemize}
There exists a finite constant \(C=C(d,\alpha,\beta,\gamma,L,\varepsilon)>0\) such that the following bound holds for every primitive law
\(P\in\mathcal M(d,\alpha,\beta,\gamma,L,\varepsilon)\) from \cref{obj:def:primitive-class}, and every \(n,m,J\in\mathbb N\) and \(h\in\mathbb R\) satisfying
\begin{itemize}
\item \textbf{(Sample sizes.)} \(n\ge2\) and \(m\ge0\).
\item \textbf{(Localization.)} \(0<h\le1/2\).
\item \textbf{(Resolution.)} \(J\ge1\).
\end{itemize}
Take expectations under \(\mathsf E_{n,m}(P)\), the product sampling law specified in \cref{obj:ass:sample-law}. Write \(N\) for the total treatment-record count, \(q\) for the number of polynomial multi-indices of total degree at most two, and \(k\) for the fine projection rank:
\[
N=n+m,\qquad q=\binom{d+2}{2},\qquad k=qJ^d.
\]
For the empirical response vector \(\widehat R_{h,J}\) and symmetrized corrected empirical matrix \(\widehat Q_{h,J}\) from \cref{obj:def:moments}, define their population expectations by
\[
\bar R=\mathbb E_{\mathsf E_{n,m}(P)}\widehat R_{h,J},
\qquad
Q_{P,h,J}=\mathbb E_{\mathsf E_{n,m}(P)}\widehat Q_{h,J}.
\]
Then
\[
\mathbb E_{\mathsf E_{n,m}(P)}
 \bigl\|\widehat R_{h,J}-\bar R\bigr\|_2^2
+
\mathbb E_{\mathsf E_{n,m}(P)}
 \bigl\|\widehat Q_{h,J}-Q_{P,h,J}\bigr\|_F^2
\le
C\left\{\frac{1}{nh^d}+\frac{k}{nNh^{2d}}\right\}.
\]
\end{lemmav}

The first term reflects labeled information within the localization cube. The second couples the two record counts and increases with the fine projection rank. Because \(\delta=h/J\), its square-root order is the rectangular sampling term in \cref{obj:thm:rectangular-upper}. The same bound covers fluctuations of the response moments and of the empirical matrix, so it supplies the sampling control needed for the guarded local fit.

\subsection{Guarded inversion and public tuning}

The decision in \cref{obj:def:estimator} uses the empirical cutoff \(g(\varepsilon)=\varepsilon(1-\varepsilon)/2\) against the population positivity margin \(\varepsilon(1-\varepsilon)\) in \cref{obj:lem:population-positivity}. This half-margin gap connects matrix fluctuations to the stability of the inversion branch. The complementary branch returns zero, and clipping confines the decision to the contrast range \([-1,1]\) supplied by \cref{obj:ass:control-interior,obj:ass:treated-interior}. These operations give the total, finite, Borel decision covered by \cref{obj:thm:rectangular-upper}; the resulting constants may depend on the fixed overlap level.

The roles of the auxiliary results are distinct. \Cref{obj:lem:local-polynomial-projection-facts} supplies the effect approximation order \(h^\gamma\), the nuisance projection orders, and bounded polynomial coefficients. \Cref{obj:lem:population-positivity} supplies population stability, while \cref{obj:lem:rectangular-sampling-bound} controls empirical departures from the population moments. Together, these ingredients underlie the four contributions in the uniform risk bound of \cref{obj:thm:rectangular-upper}: effect approximation, combined nuisance approximation, labeled sampling, and rectangular sampling. Clipping supplies the constant bound represented by the minimum with one.

The public tuning in \cref{obj:def:tuning} assigns the fine resolution according to \(S=\alpha+\beta\), the sum of nuisance smoothness orders. When \(S<\gamma\), the resolution increases as the localization scale decreases, placing nuisance approximation at the effect approximation order. When \(S\ge\gamma\), the choice \(J_*=1\) provides the required approximation control. The ceiling preserves an integer resolution while changing its order by a bounded factor.

The localization scale balances the labeled sampling requirement and the rectangular sampling requirement. Its two candidate orders correspond to the two terms in \cref{obj:def:sharp-rate}, with denominator \(\Delta=2\gamma+d+\gamma d/S\) for the unequal-count term. Taking their maximum yields the publicly tuned risk guarantee in \cref{obj:thm:rectangular-upper}. The resulting precision has a direct statistical interpretation: response labels govern the ordinary local sampling term, while the total treatment-record count reduces the sampling cost of the fine projection correction.

\subsection{The supplied-propensity upper bound}

The decision displayed in \cref{obj:thm:supplied-propensity} uses the designated propensity from \cref{obj:def:primitive-class} to form an inverse-propensity-weighted response. Under exchangeability and overlap, this response has conditional mean equal to the canonical contrast and bounded magnitude. The known uniform covariate distribution then supplies the design measure for the degree-two local polynomial average.

Polynomial reproduction makes the contrast smoothness the relevant approximation order for this decision. Its bandwidth \(n^{-1/(2\gamma+d)}\), specified in \cref{obj:thm:supplied-propensity}, balances local approximation with labeled sampling. Clipping preserves the bounded contrast range. These are the statistical ingredients of the upper-bound component of that result, which attains the order \(r_o(n)=n^{-\gamma/(2\gamma+d)}\) uniformly over the primitive class and every auxiliary-record count. This benchmark identifies the labeled-sample precision attained by the rectangular estimator when its combined approximation and sampling costs are controlled at the same order.
